% !TeX root = ../../Book1.tex
\chapter*{后记}
\addcontentsline{toc}{chapter}{后记}

本卷写到这里，最值得回想的，也许就是那几个早已熟悉的例子。\(X^3-2\) 的根并不难写出来；进一步的问题是，把这些根放在一起，究竟得到怎样的域，哪些置换能够保持根之间的代数关系，这些自同构又怎样决定中间域。求根的计算还在，问题却比求出三个数多了许多。前面逐章建立的概念，也在这样的追问中有了具体的意义。

\ProseParagraphBreak
群、环、域之间的联系，在这些例子上看得很清楚。不过，每一种对象也都有它自己的问题，值得多花一些篇幅。本卷保留了自由群的构造与有限生成交换群的分类，正是出于这样的考虑。前者允许任意指定自由生成字母的像来构造同态，后者则把一组生成元之间的关系写成整数矩阵，再用行列变换去辨认群的结构。到了后卷，有限生成交换群的分类会得到模论的解释；在这里先把整数情形做清楚，读者便有一个可以据以比较的具体结果。一般理论的意义，有时要等一个算例充分展开以后，才容易体会。

\ProseParagraphBreak
篇幅上的另一项取舍，就是把若干常用构造写得稍细一些。分式域要检查等价关系和运算；商结构要说明代表元的更换为何不影响结果；扩张上的映射则要核对原有关系是否仍然成立。这些步骤读来未必轻快，却决定了所构造的对象究竟是什么。泛性质随后给出另一种理解：一个对象能够接受或延拓哪些映射，往往比它最初采用哪套符号更能说明其特征。具体构造与这种描述各有作用；本卷希望让读者看见二者如何相互印证。

\ProseParagraphBreak
Galois 理论在本卷中占有较重的分量，这是因为它把几种看似不同的计算联系得格外紧密。对于开头提到的三次方程，先加入三次单位根，再加入一个立方根，就能把根式求解的步骤与 \(S_3\) 的正规子群链对应起来。对于 \(X^5-X-1\)，模素数后的分解则提供了辨认 Galois 群的线索。群的结构既能解释已有的根式求解步骤，也能用来证明根式求解不可能。根式求解公式的存在与否，由此成为可以证明的性质；即使没有这样的公式，方程仍然有许多值得研究的内容，例如根的算术性质及其在不同域中的行为。

\ProseParagraphBreak
无限 Galois 理论又使这种联系多了一层。有限子扩张上的自同构可以逐层考察；但把它们合在一起，就需要相容性；要从不动域恢复子群，还需要闭性。拓扑在这里有明确的代数缘由。把这一章放在卷末，是希望保留这样一次经验：已有的方法在新的对象上遇到困难时，困难本身会提示还缺少什么概念。这里建立的对应定理还没有回答绝对 Galois 群的一般结构问题；本卷只选取了几类能够具体计算的情形。

\ProseParagraphBreak
附录也作了相应的取舍。集合论与拓扑的补充用于澄清正文实际调用的知识，编码、群论及域论的专题则介绍一些可以继续研究的问题。它们各自都有远比这里丰富的理论；在有限篇幅内，所能做的，就是把选定的问题说清楚，把所给的例子算明白。读者若对其中某一题目发生兴趣，就可以沿相应的阅读提示另作系统学习。

\ProseParagraphBreak
教材的叙述容易给人一种错觉，仿佛问题一经提出，合用的定义和恰当的证明就会依次出现。实际做一道题时，往往是先有几次算不通的尝试，才看出应当研究哪个子群，或应当把哪个元素加入原来的域。写成证明时，省去了其中许多往返；学习时，却常常需要把它们补回来。一个没有做成的构造，只要能够说清它失败的缘故，也会加深对定义的理解。

\ProseParagraphBreak
这份书稿还需要在阅读和核对中继续修订。有些地方可能解释得不够，有些地方则可能说得过多；数学上的疏漏与叙述上的含混，都有待逐一改正。我希望本卷留下的，是几个能够独立算清楚的例子，一些能够自己重新证明的结果，以及由此生出的新问题。以后再遇到群、环或域时，若能想到一个值得检验的猜想，或者愿意为一个定义寻找合适的例子，那么，这些基础知识便已经开始发挥作用了。

\bigskip
\begin{flushright}
R. EverStarine

2026年9月26日
\end{flushright}
